%表紙
\newpage
\chapter{規則の改正}
\begin{enumerate}[label=第\arabic*条　]
  \item この章は、憲章の趣旨に基づき、学園規則の制定、改正及び廃止の手続に関し必要な事項を定めるものとする。
  \item 学園長又は教務主事は、学園規則の制定、改正又は廃止を発議することができる。
  \item 学園長又は教務主事は、前条の発議に当たり、必要に応じて学生の意見を募るものとする。
  \item 学園規則の制定、改正又は廃止は、発議をした者が他の学園長及び教務主事にあらかじめ通知することにより、その効力を生ずる。
  \item 学園規則を制定し、改正し、又は廃止したときは、学園長は、速やかにこれを学生に公示しなければならない。
  \item この章の実施に関し必要な事項は、学園長が別に定める。
\end{enumerate}

\bigskip
\noindent 附則

\begin{enumerate}[label=\arabic*　,leftmargin=*]
  \item この規則は、２０２６年９月１日から施行する。
  \item この規則の施行の際、現に学園の運営に従事している者は、第２章第３条の規定にかかわらず、教務主事に委嘱されたものとみなす。ただし、本人がその意思を有しない場合は、この限りでない。
\end{enumerate}
